\documentclass[manuscript,11pt,screen,nonacm]{acmart}
% Based on the unmodified official ACM manuscript sample in template/.
% STOC submission requirements override proceedings layout defaults.
\geometry{letterpaper,margin=1in,heightrounded=false,headheight=14pt,footskip=24pt}
% Keep the page folios at the same nominal size as the submission text.
\fancypagestyle{submission}{%
  \fancyhf{}
  \fancyfoot[C]{\normalfont\normalsize\thepage}
  \renewcommand{\headrulewidth}{0pt}
  \renewcommand{\footrulewidth}{0pt}}
\setcopyright{none}
\settopmatter{printacmref=false,printccs=false,printfolios=true}
\acmDOI{}
\acmISBN{}
\acmVolume{}
\acmNumber{}
\hypersetup{keeppdfinfo=true,colorlinks=true,linkcolor=blue,citecolor=blue,urlcolor=blue,
 pdftitle={A Sparse-Cardinality Curvature Bound under Persistent Adversarial Noise},
 pdfauthor={Zijian Zeng, Jiayin Cao, Weinan Dai, Rui Wang, Fei Ding},
 pdfpublisher={The authors},pdfpubtype={report},pdfpubstatus={AO},
 pdfcontactemail={1002266693@ucsiuniversity.edu.my, 2795820153@qq.com, 21725058@zju.edu.cn, 2002007@nut.edu.cn, dingfei@email.ncu.edu.cn},
 pdfcontactaddress={UCSI University; Ningbo University; Zhejiang University; Nanchang Institute of Technology; Alibaba Group},pdfmetalang={en},
 pdfsubject={Submodular maximization with bounded curvature and persistent adversarial value noise},
 pdfkeywords={submodular maximization, total curvature, persistent noise, query lower bound}}
\renewcommand{\bibliofont}{\normalfont\normalsize}
\captionsetup{font=normalsize,labelfont={normalfont,normalsize},textfont={normalfont,normalsize}}
\newcommand{\E}{\mathbb E}
\newcommand{\OPT}{\operatorname{OPT}}
\newcommand{\eps}{\varepsilon}
\newtheorem{theorem}{Theorem}
\newtheorem{lemma}[theorem]{Lemma}
\newtheorem{proposition}[theorem]{Proposition}
\newtheorem{corollary}[theorem]{Corollary}
\newcommand{\Prb}{\mathbb P}
\newcommand{\D}{\mathcal D}
\newcommand{\C}{\mathcal C}
\newcommand{\B}{\mathcal B}
\newcommand{\eq}{\mathrm{eq}}
\newcommand{\up}{\mathrm{up}}
\newcommand{\add}{\mathrm{add}}
\usepackage{needspace}
\DeclareMathOperator{\supp}{supp}
\title{A Sparse-Cardinality Curvature Bound under Persistent Adversarial Noise}
\author{Zijian Zeng}
\affiliation{\institution{UCSI University}\city{}\country{Malaysia}}
\email{1002266693@ucsiuniversity.edu.my}
\author{Jiayin Cao}
\affiliation{\institution{Ningbo University}\city{}\country{China}}
\email{2795820153@qq.com}
\author{Weinan Dai}
\affiliation{\institution{Zhejiang University}\city{}\country{China}}
\email{21725058@zju.edu.cn}
\author{Rui Wang}
\affiliation{\institution{Nanchang Institute of Technology}\city{}\country{China}}
\email{2002007@nut.edu.cn}
\author{Fei Ding}
\affiliation{\institution{Alibaba Group}\city{}\country{China}}
\email{dingfei@email.ncu.edu.cn}
% The user supplied institutions and emails; no city was supplied.
% Display all contact details at the requested minimum text size.
\authorsaddresses{\normalsize
\begin{tabular}{@{}ll@{}}
Zijian Zeng & \href{mailto:1002266693@ucsiuniversity.edu.my}{1002266693@ucsiuniversity.edu.my}\\
Jiayin Cao & \href{mailto:2795820153@qq.com}{2795820153@qq.com}\\
Weinan Dai & \href{mailto:21725058@zju.edu.cn}{21725058@zju.edu.cn}\\
Rui Wang & \href{mailto:2002007@nut.edu.cn}{2002007@nut.edu.cn}\\
Fei Ding & \href{mailto:dingfei@email.ncu.edu.cn}{dingfei@email.ncu.edu.cn}
\end{tabular}\spacefactor=3000 }
\date{}
\begin{document}
\singlespacing
\begin{abstract}
\normalsize
We study cardinality-constrained maximization of a monotone submodular
function when a fixed value oracle has adversarial multiplicative error.
The objective is the true function value, and queries may have any size.
For exact total curvature $0\le\kappa<1$ and noise $0<\eps<1$, we calibrate
a hidden-set construction to give an explicit lower bound against adaptive
randomized query algorithms. If $k=\lfloor n^\alpha\rfloor$ for fixed
$0<\alpha<1/2$ and $\eps_n k\to\infty$, every algorithm family with a fixed
polynomial query bound has fixed valid instances on which its expected
approximation is at most
$\frac{1-\eps_n}{1+\eps_n}(1-\kappa)+o(1)$.
The standard method of choosing the largest noisy singletons guarantees
$\frac{1-\eps_n}{1+\eps_n}(1-\kappa+\kappa/k)$, giving additive asymptotic
matching in this regime. For vanishing noise with $\eps_n k\to\infty$,
the limiting factor is $1-\kappa$, separated from the exact-oracle
benchmark $1-\kappa/e$ when $\kappa>0$.
\end{abstract}
\hypersetup{pageanchor=false}
\maketitle
\hypersetup{
 pdftitle={A Sparse-Cardinality Curvature Bound under Persistent Adversarial Noise},
 pdfauthor={Zijian Zeng, Jiayin Cao, Weinan Dai, Rui Wang, Fei Ding},
 pdfpublisher={The authors},pdfpubtype={report},pdfpubstatus={AO},
 pdfcontactemail={1002266693@ucsiuniversity.edu.my, 2795820153@qq.com, 21725058@zju.edu.cn, 2002007@nut.edu.cn, dingfei@email.ncu.edu.cn},
 pdfcontactaddress={UCSI University; Ningbo University; Zhejiang University; Nanchang Institute of Technology; Alibaba Group},
 pdfsubject={Submodular maximization with bounded curvature and persistent adversarial value noise},
 pdfkeywords={submodular maximization, total curvature, persistent noise, query lower bound}}
\enlargethispage{-12pt}
\thispagestyle{empty}
\clearpage
\setcounter{page}{1}
\hypersetup{pageanchor=true}
\pagestyle{submission}

\section{Model and contribution}

Curvature measures how much an element's marginal value can fall as a set
grows. With exact value access, this structure improves the approximation
factor for monotone submodular maximization to $1-\kappa/e$, up to arbitrary
accuracy slack~\cite[Theorems 6.1 and 7.1]{SVW14}. Persistent adversarial
noise can remove much of that benefit. We identify a sparse-cardinality
regime in which even unrestricted computation and adaptive randomized
queries cannot asymptotically improve a simple singleton method.

Let $V=[n]$ and $n\ge k\ge2$. The true function
$f:2^V\to\mathbb R_{\ge0}$ is normalized, monotone and submodular, and
$f(\{i\})>0$ for every $i$. Write
$f(i\mid T)=f(T\cup\{i\})-f(T)$ for $i\notin T$ and define its total curvature by
\[
 \kappa(f)=1-\min_{i\in V}
 \frac{f(i\mid V\setminus\{i\})}{f(\{i\})}.
\]
A valid oracle is a single map $g:2^V\to\mathbb R$, fixed before the
algorithm's private random choices, such that
\begin{equation}\label{eq:noise}
 (1-\eps)f(T)\le g(T)\le(1+\eps)f(T)\qquad(T\subseteq V).
\end{equation}
Repeated queries return the same value. Errors across different sets may
be arbitrarily correlated; $g$ need not be submodular or monotone.
Queries of every cardinality are permitted. The output $A$ must satisfy
$|A|\le k$, and approximation is measured by
$f(A)/\OPT(f,k)$, where $\OPT(f,k)=\max_{|T|\le k}f(T)$.
Algorithms may know $n,k,\kappa,\eps$ and use arbitrary computation.
A polynomial query family has a bound $C n^d$ for constants $C,d$ fixed
for that family. We impose no bit-complexity or numerical precision bound.

\begin{theorem}[Finite query bound]\label{thm:finite}
Let $n,k,r$ be integers with $n>k+r$ and $1\le r<k$. For
$0\le\kappa<1$ and $0<\eps<1$, set
\begin{equation}\label{eq:parameters}
 \rho=\frac{1-\eps}{1+\eps},\qquad
 a=\rho(1-\kappa),\qquad b=\rho\kappa,\qquad
 c=\frac{(1-\rho)(1-\kappa)}{1-a}.
\end{equation}
Suppose $cn>e k$. For every possibly randomized adaptive algorithm making
at most $Q$ queries on every valid input and returning a feasible set,
there is a fixed valid pair $(f_H,g_H)$, independent of its private seed,
with $\kappa(f_H)=\kappa$ and $\OPT(f_H,k)=k$, such that
\begin{equation}\label{eq:finite}
 \frac{\E[f_H(A)]}{k}
 \le a+\frac{br}{k}+\frac{(1-a)k}{n}
       +\min\left\{1,Q\left(\frac{e k}{cn}\right)^{cr}\right\}.
\end{equation}
Here $H\subseteq V$ has size $k$, and the pair is explicitly
\begin{align}
 f_H(S)&=|S\cap H|+a|S\setminus H|+b\min\{|S\setminus H|,r\},
                                                        \label{eq:family}\\
 g_H(S)&=\max\big\{(1+\eps)[a|S|+b\min\{|S|,r\}],
                         (1-\eps)f_H(S)\big\}.          \label{eq:oracle}
\end{align}
The expectation is only over the algorithm's randomness.
\end{theorem}

Almost-sure query and feasibility promises suffice. Since the proof uses
only finitely many valid inputs at each $n$, their exceptional null sets
of seeds can be discarded together.

\begin{corollary}[Sparse cardinality]\label{cor:asymptotic}
Fix $0<\alpha<1/2$ and $0\le\kappa<1$, let
$k_n=\lfloor n^\alpha\rfloor$, and suppose $0<\eps_n<1$ with
$\eps_n k_n\to\infty$. Put $\rho_n=(1-\eps_n)/(1+\eps_n)$.
For every algorithm family using at most $C n^d$ queries, there is a sequence
of fixed valid pairs of exact curvature $\kappa$ for which
\begin{equation}\label{eq:asymptotic}
 \frac{\E[f_n(A_n)]}{\OPT(f_n,k_n)}
 \le \rho_n(1-\kappa)+o(1).
\end{equation}
On every valid input, selecting the $k_n$ largest noisy singleton values
guarantees $\rho_n(1-\kappa+\kappa/k_n)$.
The bounds match additively; they match relatively when $\rho_n$ is bounded
away from zero.
\end{corollary}

The instances may depend on the algorithm family, but each oracle is fixed
independently of its private seed. The statement quantifies separately over
each fixed polynomial query bound. It does not optimize over unbounded
polynomial degrees at a fixed $n$. For example, $\kappa=1/2$ and
$\eps=1/3$ yield the matching limiting factor $1/4$.
If $\eps_n\to0$ while $\eps_n k_n\to\infty$, the factor tends to
$1-\kappa$; for $\kappa>0$ this is separated from the exact-oracle benchmark.

\paragraph{Relation to prior work.}
Horel and Singer~\cite[Section 2.1 and Proposition 8]{HorelSinger16}
study an approximately submodular objective $F$ satisfying the same
pointwise sandwich as~\eqref{eq:noise}. They supply both the hidden
capped-rank pattern and the largest-noisy-singleton method. Their
curvature factor $(1-\kappa)\rho^2$ is stated for the approximate
objective $F$. Directly analyzing the same singleton method against
true $f$ gives one factor of $\rho$, as proved below. That normalization
and its $\kappa/k$ refinement are elementary baselines.
Sviridenko, Vondr\'ak and Ward~\cite[Section 7]{SVW14} use a uniform
modular shift to tune curvature in exact-oracle lower bounds.
Hassidim and Singer~\cite[Theorem 8 and Section 5]{HassidimSinger17}
already establish adversarial-noise hardness for randomized algorithms.
Their displayed hardness regime does not prescribe the exact curvature
below one used here. Persistence also appears in stochastic-noise
algorithms: Bhawalkar, Cai, Feng, Liaw and Lin~\cite[Section 2]{BCFLL25}
require unbiased noise and independence across distinct sets, assumptions
absent from~\eqref{eq:noise}.

The calibration in~\eqref{eq:parameters} and finite bound~\eqref{eq:finite}
yield~\eqref{eq:asymptotic}
for every sequence $\eps_n k_n\to\infty$ in the stated regime.
All-query-size leakage control is needed because queries are unrestricted.
The construction and transcript strategy build on the preceding work.

\section{Construction and leakage}

\subsection{Global oracle validity and exact curvature}

Use the parameters of Theorem~\ref{thm:finite}. For $s=|S|$,
$h=|S\cap H|$ and $t=s-h$, define the baseline
\begin{equation}\label{eq:baseline}
 F_0(S)=a s+b\min\{s,r\},\qquad g_0(S)=(1+\eps)F_0(S).
\end{equation}
All formulas define values on every subset before an algorithm is run.
The marginal of an element of $H$ in $f_H$ is always one. Outside $H$,
it is $a+b=\rho$ while fewer than $r$ outside elements are present,
and $a$ thereafter. These marginals are positive and diminishing, so
$f_H$ is normalized, monotone and submodular. Since $n-k>r$, the
full-set marginal of an outside element is $a$.

\begin{table}[htbp]
\caption{Calibration of the hard family. The noisy intervals of the two
singleton types meet at $1-\eps$; both can receive that same answer.}
\label{tab:calibration}
\centering
\begin{tabular}{lccc}
\toprule
Element type & True singleton & Full-set marginal & Marginal/singleton\\
\midrule
Inside $H$ & $1$ & $1$ & $1$\\
Outside $H$ & $\rho$ & $\rho(1-\kappa)$ & $1-\kappa$\\
\bottomrule
\end{tabular}
\end{table}

Table~\ref{tab:calibration} proves that $\kappa(f_H)=\kappa$, including
$\kappa=0$. Also $f_H(S)\le h+(a+b)t\le s$ and $f_H(H)=k$,
so the feasible optimum is exactly $k$.
Replacing $s$ by $s-h$ decreases its capped value by between zero and $h$.
It follows that
\begin{equation}\label{eq:difference}
 (1-\rho)h\le f_H(S)-F_0(S)\le(1-a)h.
\end{equation}
Thus $F_0(S)\le f_H(S)$. Both terms in the maximum~\eqref{eq:oracle}
are at most $(1+\eps)f_H(S)$, and the second term supplies the lower
endpoint in~\eqref{eq:noise}. This verifies the multiplicative promise
globally, including unqueried sets, large sets and the empty set.

The baseline pair $(F_0,g_0)$ is also a valid input of exact curvature
$\kappa$: its singleton value is $\rho$ and its full-set marginal is $a$.
We may therefore use the algorithm's query and output promises on this
baseline. Finally,
\begin{equation}\label{eq:c}
 0<c\le1,\qquad c\ge(1-\kappa)\eps.
\end{equation}
Indeed the denominator defining $c$ is positive, and subtracting its
numerator from that denominator gives $\kappa$. The second bound follows
from $1-\rho=2\eps/(1+\eps)\ge\eps$ and $1-a\le1$.

\needspace{8\baselineskip}
The unequal singleton values are essential to this calibration. In contrast,
the particular uniformly shifted family
\[
 U_H(S)=(1-\kappa)|S|+
             \kappa\bigl(|S\cap H|+\min\{|S\setminus H|,r\}\bigr)
\]
has optimum $k$ and gives every $k$-set value at least $(1-\kappa)k$.
It cannot witness the smaller fixed-noise factor $\rho(1-\kappa)$.
The calibration in Table~\ref{tab:calibration} retains exact curvature
while lowering the outside marginal to that product.

\subsection{A tail bound valid for all query sizes}

Let $H$ be uniform among the $k$-subsets of $V$.
If $s\le r$, then $f_H(S)=\rho s+(1-\rho)h\le s$ and
$g_0(S)=(1-\eps)s$, so $g_H(S)=g_0(S)$ for every $H$.
For $s>r$, the upper bound in~\eqref{eq:difference} implies equality
of the oracle answers whenever
\begin{equation}\label{eq:theta}
 h\le\theta(s):=
 \frac{(1-\rho)[(1-\kappa)s+\kappa r]}{1-a}\ge cs.
\end{equation}
To see this, equality is guaranteed by
$(1-\eps)[F_0(S)+(1-a)h]\le(1+\eps)F_0(S)$;
solving this inequality gives~\eqref{eq:theta}.

Fix an $s$-set $S$ with $s>r$ and let $j=\lceil cs\rceil$.
By~\eqref{eq:c}, $1\le j\le s$. A revealing answer requires
$|S\cap H|>cs$, hence lies in the enlarged event
$|S\cap H|\ge j$, even when $cs$ is an integer. If $j>k$ that event
is empty. Otherwise a union bound over the $j$-subsets of $S$ gives
\begin{align}
 \Prb[g_H(S)\ne g_0(S)]
 &\le \binom{s}{j}\frac{(k)_j}{(n)_j}
 \le \binom{s}{j}\left(\frac{k}{n}\right)^j \notag\\
 &\le\left(\frac{e s k}{j n}\right)^j
 \le\left(\frac{e k}{cn}\right)^j
 \le\left(\frac{e k}{cn}\right)^{cr}.             \label{eq:tail}
\end{align}
Here $(m)_j=m(m-1)\cdots(m-j+1)$ and
$(k-i)/(n-i)\le k/n$. The standard estimate
$\binom{s}{j}\le(e s/j)^j$ follows from $j!\ge(j/e)^j$.
The last inequality uses $cn>e k$ and $j\ge cs>cr$.
No independence among membership indicators is used.
Together with the nonrevealing case $s\le r$, this controls every query size.

\section{Adaptive queries and fixed hard instances}

\begin{proof}[Proof of Theorem~\ref{thm:finite}]
The construction already verifies the structural promises and optimum.
Fix a private seed $\omega$ and run the algorithm on $g_0$.
Let $S_1(\omega),\ldots,S_m(\omega)$ be its baseline queries, where
$m\le Q$, and let $A_0(\omega)$ be its output. All these sets are
independent of $H$. Let $B$ be the event that at least one baseline
query has different answers under $g_H$ and $g_0$.
By~\eqref{eq:tail} and a union bound,
\begin{equation}\label{eq:bad}
 \Prb_H(B\mid\omega)\le
 \eta:=\min\left\{1,Q\left(\frac{e k}{cn}\right)^{cr}\right\}.
\end{equation}
When $B$ does not occur, induction through the successive queries shows
that the run on $g_H$ has the same transcript, stopping time and output
as the baseline run. The induction uses independence from $H$ only for
the baseline queries; it makes no assertion about queries after a divergence.
Every feasible output has true value at most $k$. Pointwise, therefore,
\[
 f_H(A^{g_H}_\omega)\le f_H(A_0(\omega))+k\mathbf 1_B.
\]
The baseline output need not have been queried. Nevertheless, with
$m_0=|A_0|\le k$, symmetry gives $\E_H|A_0\cap H|=m_0 k/n$.
Using~\eqref{eq:family},
\begin{align*}
 \E_H[f_H(A_0)]
 &=a m_0+(1-a)\E_H|A_0\cap H|
                  +b\E_H\min\{|A_0\setminus H|,r\}\\
 &\le ak+(1-a)k^2/n+br.
\end{align*}
Average over $H$ and then over $\omega$, and apply~\eqref{eq:bad}.
This yields~\eqref{eq:finite} on the joint average.
Since the collection of $H$ is finite, at least one single $H$ has
private-seed expected ratio at most that average. Its oracle is the fixed
map~\eqref{eq:oracle}, independent of the seed, as required.
\end{proof}

\section{Matching algorithm and asymptotic regime}

\begin{proposition}[Standard singleton method]\label{prop:singletons}
For any $n\ge k\ge2$, $0\le\eps<1$, and a valid pair $(f,g)$
with total curvature $\kappa$, selecting the $k$ largest $g(\{i\})$
values returns $S$ with
\begin{equation}\label{eq:singleton}
 f(S)\ge\rho(1-\kappa+\kappa/k)\OPT(f,k).
\end{equation}
\end{proposition}

\begin{proof}
Let $O$ be an optimal $k$-set, extending a smaller optimizer if necessary.
Singleton ordering and submodularity imply
\begin{align*}
 (1+\eps)\sum_{i\in S}f(\{i\})
 &\ge\sum_{i\in S}g(\{i\})
 \ge\sum_{i\in O}g(\{i\})\\
 &\ge(1-\eps)\sum_{i\in O}f(\{i\})
 \ge(1-\eps)f(O).
\end{align*}
Enumerate $S$ starting with its largest true singleton. Diminishing
marginals and curvature give every later marginal at least
$(1-\kappa)f(\{i\})$. Hence
\[
 f(S)\ge(1-\kappa)\sum_{i\in S}f(\{i\})
                  +\kappa\max_{i\in S}f(\{i\})
 \ge(1-\kappa+\kappa/k)\sum_{i\in S}f(\{i\}).
\]
Combining the inequalities proves the claim. The true singleton order is
used only in the analysis; the algorithm queries and sorts noisy values.
This is the direct true-objective analysis of the method in
Horel--Singer~\cite[Proposition 8]{HorelSinger16}.
\end{proof}

\begin{proof}[Proof of Corollary~\ref{cor:asymptotic}]
Suppress subscripts and set $u=\eps k\to\infty$,
$\delta=1-\kappa>0$, and $r=\lceil\sqrt{k/\eps}\rceil$.
Then
\[
 \frac r k\le\frac1{\sqrt u}+\frac1k\to0.
\]
Eventually $1\le r<k$ and $n>k+r$. By~\eqref{eq:c},
\begin{equation}\label{eq:growth}
 cr\ge\delta\sqrt u\to\infty,\qquad
 \frac{cn}{k}\ge\delta u\frac{n}{k^2}
                      \ge\delta u n^{1-2\alpha}.
\end{equation}
Put $\gamma=(1-2\alpha)/2>0$. Eventually
$e k/(cn)\le n^{-\gamma}<1$, so Theorem~\ref{thm:finite} applies.
For the fixed query bound $Q\le C n^d$, the disagreement term is at most
\[
 C n^{d-\gamma cr}=n^{-\omega(1)}.
\]
The remaining error terms satisfy
\[
 \frac{br}{k}\le\kappa\left(\frac1{\sqrt u}+\frac1k\right)\to0,
 \qquad \frac{(1-a)k}{n}\le\frac{k}{n}\to0.
\]
Applying the finite theorem separately at each sufficiently large $n$
gives~\eqref{eq:asymptotic}. Proposition~\ref{prop:singletons} supplies
the deterministic lower guarantee. Its gap above $\rho_n(1-\kappa)$
is $\rho_n\kappa/k_n=o(1)$, giving additive matching. When
$\inf_n\rho_n>0$, the central factor is bounded below by a positive
constant, so the additive remainders also give relative matching.
\end{proof}

\section{Exact finite checks and boundaries}

The same transcript proof has a finite alternative that avoids the
condition $cn>e k$. Retain $n>k+r$, $1\le r<k$ and the other
parameters of Theorem~\ref{thm:finite}, and define
\begin{align}
 p&=\max_{0\le s\le n}
       \Prb_H[g_H(S_s)\ne g_0(S_s)],\label{eq:p}\\
 \beta&=a+\frac{(1-a)k}{n}
              +\frac b k\E[\min\{k-X,r\}],\qquad
 X\sim\operatorname{Hypergeom}(n,k,k),\label{eq:beta}
\end{align}
where $S_s$ is any fixed $s$-set. Symmetry makes~\eqref{eq:p}
independent of its choice.

\begin{lemma}[Exact-tail alternative]\label{lem:exact}
For every $Q$-query algorithm as in Theorem~\ref{thm:finite}, some
fixed $H$ satisfies
\begin{equation}\label{eq:exact}
 \frac{\E[f_H(A)]}{k}\le\beta+\min\{1,Qp\}.
\end{equation}
\end{lemma}
\begin{proof}
The baseline union bound is at most $\min\{1,Qp\}$. Extend any
baseline output of size below $k$ to a $k$-set and use monotonicity.
For every fixed $k$-set, its expected true value divided by $k$ is
exactly~\eqref{eq:beta}, by the hypergeometric law of its intersection
with $H$. The same coupling and finite averaging prove~\eqref{eq:exact}.
\end{proof}

All probabilities in~\eqref{eq:p} can be computed by summing
\[
 \frac{\binom{s}{h}\binom{n-s}{k-h}}{\binom{n}{k}}
\]
over feasible $h$ for which the formulas~\eqref{eq:baseline}
and~\eqref{eq:oracle} differ. For $n=6,k=3,r=2$, $\eps=1/2$ and
$\kappa=1/3$, this gives $p=1/20$ and $\beta=359/540$.
Exhaustive optimization of deterministic adaptive decision trees of
depth zero, one and two gives respective best uniform-prior expected
ratios $359/540$, $7/10$ and $13/18$, each satisfying~\eqref{eq:exact}.
A randomized policy is a mixture of deterministic policies, so it cannot
improve those prior optima. The portable supplementary script reproduces
these checks and verifies additional all-subset and count-based instances.
These finite observations supplement the analytic proofs.

The curvature-one endpoint, the transition with bounded $\eps k$,
and dense cardinality remain open directions for a fuller joint-parameter
characterization. The case $k=n$ provides a boundary condition: returning
$V$ is optimal. Further questions include exact finite minimax factors
and high-probability guarantees on fixed adversarial inputs.

\paragraph{Computational and AI assistance.}
Codex assisted with the construction, proof development, exact finite checks,
literature comparison and drafting. A separate model-context pass checked
the mathematical arguments.

\bibliographystyle{alpha}
\bibliography{references}
\clearpage
\onehalfspacing
\end{document}
